% ============================================================
%  R. Nielsen, "Forelæsninger over »Philosophisk Propædeutik«"
%  fra Universitetsaaret 1860–61
%  Gyldendalske Boghandling, Kjøbenhavn, 1862
%
%  Excerpt: pp. 326–328
%  (PDF pages 338–340 of the Google Books scan)
%
%  These three pages contain the main philosophical engagement
%  with Darwin in the 1862 lectures. They fall at the end of
%  §II.C ("Philosophiens Forhold til Physiologien: Livsproblemet"),
%  immediately after Nielsen completes his positive argument for
%  the dialectical reality of species-ideas (Artsideer), and just
%  before the section on animal life (β: Dyrelivet) begins.
%
%  Darwin is known to Nielsen through H. G. Bronn's German
%  translation, Ueber die Entstehung der Arten im Thier- und
%  Pflanzenreich (Stuttgart 1860), which had appeared just two
%  years before these lectures. (Jacobsen's Danish translation
%  would not appear until 1871–72.)
%
%  Note on p. 327: In the original, Nielsen's footnote reference
%  (*) on p. 326 introduces a note so long that it fills all of
%  p. 327, leaving only three lines of body text at the top.
%  Rather than a \footnote{} command, which LaTeX would struggle
%  to accommodate, the note is rendered here as an indented block
%  quotation, preceded by an editorial marker [*] in the body text.
%
%  Transcription from the Google Books digitisation (Royal
%  Library Copenhagen copy). OCR errors corrected by
%  inspection of the scan images. Orthography follows the
%  original throughout.
% ============================================================
\documentclass[12pt,a4paper]{article}
\usepackage[utf8]{inputenc}
\usepackage[T1]{fontenc}
\usepackage[danish]{babel}
\usepackage{microtype}
\usepackage{setspace}
\usepackage[margin=1.2in]{geometry}
\usepackage{fancyhdr}
\setlength{\headheight}{14.5pt}
\usepackage{hyperref}

\hypersetup{
  pdftitle={Forelæsninger over »Philosophisk Propædeutik« (uddrag)},
  pdfauthor={R. Nielsen},
  hidelinks
}

\pagestyle{fancy}
\fancyhf{}
\rhead{R.\ Nielsen, \emph{Phil.\ Prop.}\ (1862), s.\ 326--328}
\cfoot{\thepage}

\setstretch{1.3}

\title{\textbf{Forelæsninger over »Philosophisk Propædeutik«}\\[0.4em]
\large fra Universitetsaaret 1860--61\\[0.3em]
\normalsize Uddrag: s.\ 326--328\\
(om Arternes Tilblivelse og Darwins Theori)}
\author{R.\ Nielsen}
\date{Gyldendalske Boghandling (F.\ Hegel), Kjøbenhavn, 1862}

\begin{document}
\maketitle
\thispagestyle{fancy}

\bigskip

% ------------------------------------------------------------------
%  S. 326–328  [§II.C: Philosophiens Forhold til Physiologien,
%  sub-section on plant life and the classification of organisms]
%
%  Context: Nielsen has just completed his positive argument for the
%  dialectical reality of species-ideas (Artsideer) and the law of
%  concretion (Concretionsloven), according to which genus is more
%  concrete than species, family more concrete than genus, and so on.
%  A key corollary is that there is no a priori developmental law
%  determining how many species belong to each genus: philosophical
%  cognition proceeds from facts upward to concrete ideas, not
%  downward from abstract ideas to facts. Darwin is introduced as
%  the strongest contemporary challenge to the reality of types.
%
%  The passage on p. 326 begins mid-sentence from p. 325; the text
%  on p. 328 breaks off as the sub-section β (Dyrelivet) begins.
%
%  The note marked [*] corresponds to Nielsen's footnote reference
%  on p. 326. In the original it runs so long — quoting Bronn's
%  Schlußwort at length — that it occupies all of p. 327, leaving
%  only three lines of body text at the top of that page.
% ------------------------------------------------------------------

\noindent\textit{[S.\ 326, fortsat fra s.\ 325:]}

\medskip

\noindent gjennem alle sine Arter til Familien o.\ s.\ v., kan
dialektisk bestemmes ved Concretionsloven; derimod gives der ikke
nogen Udviklingslov, hvorved det a priori kan bestemmes,
hvormange og hvilke Arter der nødvendig maae høre til hver Slægt,
eller hvormange og hvilke Slægter der oprindelig høre til hver
Familie: den philosophiske Erkjendelse gaaer ikke gjennem
abstracte Ideer til virkelige Kjendsgjerninger, men omvendt
gjennem de virkelige Kjendsgjerninger til de concrete Ideer.
(Jvfr. \emph{Phil.\ Prop.}\ S.\ 158--178).

Dersom de Forsøg, man i den nyere Tid har gjort paa at forklare
Arternes Tilblivelse af Varieteter, som i Tidernes Løb ere blevne
constante, skulle forstaaes saaledes, at det egentlig er
Omstændighederne og de ydre Betingelser, der bestemme de organiske
Typer: da maa en saa principløs og ideefornegtende Opfattelse af
Livsvirksomheden naturligviis møde Modsigelse i Philosophien. Men
det bør i denne Sammenhæng dog heller ikke glemmes, at Ideernes
Realitet er uendelig dialektisk, og at en dialektisk Realitet ikke
lader sig hypostasere i dogmatisk positiv Umiddelbarhed. Forholder
det sig virkelig saaledes, at den organiske Selvhed oprindelig
forudsætter den elementære Andethed, saa følger jo ligefrem heraf,
at Livsudviklingen ogsaa maa medføre en mangfoldig Vexelvirkning
af det Typiske og det Elementære. Varieteter, Aborter,
Forskydninger og andre Abnormiteter vise tydelig nok, hvilken
Indflydelse de elementære Betingelser maae have paa Paradigmets
Udførelse. [*]

\medskip
\noindent\textit{[*] Anmærkning (s.\ 326--327 i originalen):
Hvad de herhen hørende Undersøgelser angaaer, henvise vi til
Charles Darwin \emph{über die Entstehung der Arten im Thier- und
Pflanzenreich.} Aus dem Englischen übersetzt und mit Anmerkungen
versehen von Dr.\ H.\ G.\ Bronn. Stuttgart 1860. --- Det følgende
er citeret fra \emph{,,Schlußwort des Übersetzers''}, S.\ 508--9:}

\begin{quote}
\itshape
,,Es sind die Existenz-Bedingungen, welchen sich die Organismen
anpassen müssen, und welche eine so große Rolle in diesem Buche
spielen. Sie sind theils organische und theils unorganische, und
die ersten nach Hrn.\ Darwin weitaus die zahlreichsten und
wichtigsten und daher auch an und für sich geeignet, die
manchfaltigsten Folgen zu veranlassen. Doch eben diese organischen
Prinzipien haben für Darwin wieder den großen Vortheil, daß,
indem er sich auf ihre Manchfaltigkeit und auf den Kampf ums
Daseyn überhaupt beruft, er der Nothwendigkeit überhoben ist,
Rechenschaft von ihrer Wirkungs-Weise im Einzelnen zu geben und
nachzuweisen, welche spezielle Folgen diese oder jene spezielle
organische Bedingung auf die Struktur und Entwickelung der ihrem
Einfluß unterliegenden Organismen überhaupt, oder auf einzelne
insbesondere ausübt. Hr.\ Darwin beruft sich auf jeder Seite
darauf, daß nur solche Abänderungen Aussicht auf Erhaltung haben,
welche dem Individuum und somit der künftigen Spezies nützlich
sind; und theoretisch muß man zugestehen, daß, woferne es eine
natürliche Züchtung gebe, die Sache sich nicht anders verhalten
könne. Aber wir müssen gestehen, doch in fast allen unseren aus
angeblich innern Ursachen hervorgegangenen Varietäten gar nicht
finden zu können, worin denn der Nutzen ihrer Abänderung
bestehe\ldots. Warum bekommt z.\ B.\ in diesem Kampfe ums Daseyn
eine Pflanzen-Art ovale statt lanzettlicher und die andre
lanzettliche statt ovaler Blätter? warum die eine einen
Dolden-artigen und die andre einen Rispen-förmigen Blüthenstand?
warum die eine fünf und die andre vier Staubgefässe, die eine
eine geschlossene und die andre eine weit geöffnete Blüthe? Wozu
nützt der einen Dieß und der andern das Gegentheil? Warum
bewirken die organischen Bedingungen Dieß? Mit welchen Mitteln
fangen sie es an? und wie müssen sie beschaffen seyn, um es zu
können? Und wie kann die eine Art der andern dadurch überlegen
werden? Wir gestehen, keinen Zusammenhang zwischen diesen
Erscheinungen zu erkennen, und Hr.\ Darwin würde uns antworten,
daß es möglicher Weise so oder so zugehen könne. Wir gestehen
ferner, uns vergeblich um positive Beweise oder auch nur Belege
in dieser Beziehung umgesehen zu haben, manche spezielle Fälle
eigenthümlicher Art etwa ausgenommen; denn wir sind weit entfernt
davon, allen solchen Einfluß überhaupt läugnen zu wollen.''
\end{quote}

\noindent\textit{[S.\ 327: Brødteksten begynder øverst på siden
med tre linjer umiddelbart efter anmærkningens afslutning:]}

\medskip

\noindent ,,Den rene Fornufts Kritik'' forvandler Ideen til et
subjectivt Ubetinget; den dogmatiske Positivisme hypostaserer
Ideen til et for sig bestaaende Urvæsen. Et saadant positivt

\medskip
\noindent\textit{[S.\ 328, fortsat fra s.\ 327:]}

\medskip

\noindent Urvæsen erklæres nu rigtignok for et indvortes,
ubilledligt og usynligt Væsen, men bliver just, idet Positiviteten
slaaes fast, ved sin dogmatiske Umiddelbarhed til et udvortes,
billedligt og anskueligt Væsen. Den dybe Skuen, den umiddelbare
Idee-Skuen, er i Almindelighed en mystisk-dogmatisk Skuen af et
Negativt-Positivt, et Ubilledligt-Billedligt; det Dybe kjendes da
især derpaa, at den Skuende ikke mærker Modsigelsen. Skal den
umiddelbare Idee-Skuen faae videnskabelig Betydning, kan det kun
skee derved, at Kritiken ikke blot paaviser, men ogsaa
gjennemreflecterer den fra alle Sider frembrydende Modsigelse.
Modsigelsen er, at det i Naturen Umiddelbare er saa reflecteret;
Modsigelsen er, at det saaledes Reflecterede er saa umiddelbart.
Saalidt som den exacte Uendelighedsanalyse er færdig dermed, at
hver endelig Størrelse opfattes som et Integral af uendelig smaa
Størrelser, ligesaa lidt er den dialektiske Idee-Analyse afsluttet
dermed, at det Umiddelbare reflecterer sig i det Reflecterede; thi
de articulerende Mellembestemmelser, som heraf udvikle sig, ere
utallige.

Ideeriget er i Forhold til den ydre, deri begrundede Planteverden
et System af Reflexionsexistenser, indholdstrige
Existenserindringer; Ideernes Reflexion i Individualformernes
reflecterede Umiddelbarhed er, som Lysets Klarhed, en forklaret
Umiddelbarhed. I denne Klarhed sées Exemplarexistenserne gjennem
deres Individualideer reflecterede i Artens Idee, Artsideerne med
deres Individualideer reflecterede i Slægtens Idee, Slægtsideerne
med deres Artsideer i Familiens, gjennem Familiens i Ordenens,
gjennem Ordenens i Klassens, gjennem Klassens i Rækkens, gjennem
Rækkens concrete, indholdstrige Idee i Rigets Grundidee, i Dybet.

\medskip
\noindent $\beta$) \textbf{Dyrelivet.}

Som organisk Natur er Ideen Liv, og den levende Existens,
væsentlig seet, en Idee-Existens; men Idee-Existensen

\medskip
\noindent\textit{[Fortsat på s.\ 329, ikke medtaget her.]}

\end{document}
